\section{Reusable Lessons}
\label{sec:lessons}

The eight diagnostics converge on a small set of portable lessons. We state them
as design rules, each tied to the apparatus failure that motivated it.

\begin{enumerate}
  \item \textbf{Every criterion must carry its own counter-evidence.} A criterion
  that cannot turn red is decoration. Across the eight, the failure modes that
  produced fake-green verdicts were a quantization floor (NSP slope), a
  mis-aligned evaluation arm (multi-view), a template shortcut (DoM probe), a
  cancellation artifact from a correlation that was true for all late steps
  (termination lag), a ``nan'' returned by a non-central-$t$ CDF for large
  non-centrality (workflow power), a reversed criterion direction (1.9x), and a
  slow run that disagreed with a published table because a caliber parameter was
  wrong (AROMA). Each was caught only because we injected a known violation and
  required a red verdict.

  \item \textbf{Fix the criterion before the run, then let it fail.} The NSP
  three-arm is the clearest case: the pre-registered ``lowest slope'' criterion
  failed, and we reported that rather than swapping in the secondary finding as
  the headline. The secondary finding (lowest absolute error) is stated as a
  secondary finding, with the reason the slope criterion is unreliable on stable
  dynamics.

  \item \textbf{Diagnose apparatus defects before changing them.} Each of the
  fixes above was preceded by a concrete diagnosis (the quantization floor, the
  mis-aligned arm, the nan branch), not by a tweak that made the number look
  better. Where we changed a prior from 0.75 to 0.95 in the dynamics, it was a
  physical argument---a strongly stable system attenuates error so much that the
  arms become indistinguishable---made before observing the result, not after.

  \item \textbf{``Substrate unavailable'' is not the same as ``no conclusion''.}
  Power analysis converts ``cannot measure'' into ``MDE = 0.234, R = 16 would
  suffice''; arithmetic self-check converts ``not auditable'' into ``the feasible
  interval excludes the headline''; null geometry converts ``method paradox''
  into ``r $\ll$ d arithmetic''; stratified re-computation converts a
  ``real-robot experiment'' mean into a distributional question answerable from
  a published table. The four unavailable-substrate cards produced four citable
  quantitative statements, not four shrugs.

  \item \textbf{Conclusion scope is fixed by the apparatus, not the model.}
  This is the paper's central claim. The same nominal quantity (0.66) means
  three different things across the embodied literature (a sim-to-real
  reproduction rate, a parts Jaccard index, and an AUC with p = 0.014); 3.1\% and
  $-$7.9\% belong to two different models and cannot be merged into an interval;
  1.9x and 9.3x are different calibers and cannot be ranked. Every number must
  travel with its ruler.

  \item \textbf{Negative and non-result verdicts are first-class.} Four of the
  eight diagnostics return a failed criterion, a mixed result, a cannot-judge,
  or a not-identifiable verdict. We count these as results because they bound
  what can be claimed, and because silently converting them into positive
  findings is exactly the narrative-to-evidence error this paper targets.
\end{enumerate}

Together these rules operationalize the three engineering principles that emerge
from the six directions: \emph{verifiable state prediction} (supervision anchored
to physical or experimental quantities), \emph{auditable failure modes}
(evaluation decomposed to specific failures with variance or intervals), and
\emph{offline generation plus verification} (execution decoupled from checking so
failures can be replayed). Each principle is an instance of the same discipline
the eight diagnostics enact: make the claim checkable, make the check falsifiable,
and state what the check does not cover.